\section{Introduction}
\label{sec:geometric-results}
A near-onset alternating collision record has a stationary action with
two free endpoints, irrespective of its length. Its probability,
however, becomes exponentially small. The question addressed here is
whether the nonlinear physical law admits a relative limit on a
contact neighborhood that does not shrink with the number of flights,
and whether the resulting invariant can be recovered from actual
collision reports. We answer this question for smooth periodic
dispersing billiards without symmetry or a finite-horizon assumption.

There are three different levels of information. The leading selected
count and endpoint-metric amplitudes depend on the gap, free area and
two contact curvatures, but do not determine higher contact jets.
Nonlinear short-record probabilities contain more information: for two
independent even contacts, the two oriented two-flight germs determine
every finite pair of even jets. The long-bridge limit has a different
role. It separates the nonlinear physical amplitude into two half-line
factors and identifies a complete symmetrized energy profile at each
contact, without assuming evenness. Finite-jet determination, relative
factorization, and smooth function-valued determination are not
interchangeable conclusions.

The relative estimate requires more than control of the stationary
length. The endpoint Hessian stays positive, whereas its mixed
endpoint derivative, which is the physical flux, decays exponentially.
An absolute error divided by that derivative need not be small. We
normalize the exact cofactor identity before comparing the finite
interior Hessian with its two half-line limits. The nonlinear Hessian
perturbations are summable in trace norm near the two ends. This gives
relative determinant control without a factor proportional to the
number of interior impacts. A common Morse domain then passes the
estimate through residual-time integration, including all fixed mixed
geometric and offset derivatives down to the onset.

The resulting same-type laws, denoted $\mathcal F_{00}$ and
$\mathcal F_{11}$, determine the normalized energy profiles
$\mathcal V_0,\mathcal V_1$ on a common closed collar. Each profile
is the pushforward of a half-line determinant-weighted endpoint measure
by its excess action, with the two transverse branches summed.
The mixed-type law is determined by these profiles and satisfies a
Volterra convolution identity. Uniqueness is proved as a Volterra
integral equation, not by matching Taylor series of smooth functions.
This addresses the inverse problem for the whole invariant rather
than only its formal jets. Separately observed odd records provide a
compatibility test for the invariant reconstructed from the even ones.

\subsection{Geometric information at finite and infinite flight number}
For the even-contact problem the gap and the two labelled curvatures
are supplied individually. The contacts need neither be congruent nor
have distinct curvatures. Two block-triangular maps recover both finite
jet vectors: Theorem~\ref{thm:v12-two-contact} uses the limiting laws,
while Theorem~\ref{thm:v13-two-flight} uses two flights directly.
If $z=(1+g\kappa_0)(1+g\kappa_1)-1$, the separating factor for the
latter is
\[
 (1+2z)^m-(1+2mz)=\sum_{r=2}^m\binom mr(2z)^r>0,
                         \qquad m\ge2.
\]
For the limiting map it is
$m\tanh((m-1)\gamma)-(m-1)\tanh(m\gamma)>0$.
The two nonvanishing factors arise from different actions and
amplitudes. Corollary~\ref{cor:v13-jet-coordinates} identifies the
resulting finite-jet coordinate changes, without identifying the
finite and limiting law functions or their noisy experiments.

Theorem~\ref{thm:v12-realization} realizes independent contact jets in
analytic periodic tables with the entire selected leading hierarchy
fixed. Its analytic coordinate family therefore gives physical,
not merely formal, alternatives invisible to those leading data.
Theorem~\ref{thm:v13-two-flight-observation} proves that $2(M-1)$
positive windows at the single flight number $j_0=2$ give local
coordinates for jets through order $2M$ in each such fixed family.
The resulting confidence cost is $O(\varepsilon^{-2}\log(1/\eta))$
and the squared minimax risk is of order $N^{-1}$. These are the
regular parametric orders after physical local invertibility is
proved. Their role is to complete this finite-dimensional physical
experiment; they are not the explanation for the long-bridge theorem.
The offsets and conditioning constants may depend on $M$.

Analytic continuation then determines the participating analytic
boundaries from either complete pair of germs in the even class.
The continuation step itself is classical. The geometric content here
is the explicit independent-contact inverse and its physical image,
including equal curvatures. The function-valued results return to
arbitrary smooth contacts and determine their symmetrized energy
profiles, without an evenness restriction or a supplied finite-dimensional
family. Thus the general invariant, rather than a supposed necessity
of long records for finite jets, is the principal inverse object.

\subsection{Observation of the complete invariant}
Write the integrated flux as
$H_{j,b}(d)=2A\sinh(j\gamma)\Pr(E_{j,b}(d))$.
On differences of limiting physical fluxes, a second derivative
followed by a weighted Abel half derivative gives a two-sided Lipschitz
coordinate for the profiles (Theorem~\ref{thm:v11-abel-inverse}).
The affine nullspace of this seminorm vanishes on the physical
flux class. Its exact monomial linearization has multiplier of order
$n^{-5/2}$; this is why a uniform error in the raw probabilities
cannot simply be substituted for the differentiated inverse norm.

For profiles bounded in $C^{m-1,1}$, $m\ge4$, the integrated flux
has two additional derivatives. Positive-node reconstruction exploits
that gain while treating smooth bridge error separately from scalar
sampling noise. Theorem~\ref{thm:v12-acquisition} gives uniform
accuracy $\varepsilon$, with confidence $1-\eta$, using at most
\[
 C\varepsilon^{-\left(2+6/(m-1/2)+\gamma/|\log\tau|\right)}
                       \log\frac{C}{\eta\varepsilon}
\]
independent binary preparations. Every failure is charged; neither
values at zero nor derivatives are observed. This is a sufficient
bound for the specified procedure, not a minimax assertion. The
number $\tau$ is a supplied relative convergence certificate.

Theorem~\ref{thm:v12-self-calibrated} removes the actual gap, area
and multiplier by a separately charged smooth-class pilot using one-
and two-flight bits. A gap error translates the integrated flux before
regularization, and hence has no spurious division by a small offset.
The pilot does not increase the uniform sufficient exponent, with
$\gamma$ replaced by its supplied upper bound. Labels, coarse boxes,
regularity and convergence certificates remain part of the experiment.
The theorem does not recover separate contact curvatures from the
multiplier. The auxiliary endpoint, count-only, coalescence and
smooth-remainder experiments retain their own hypotheses and full
proofs in the appendix.

\subsection{Geometric setting and observations}
Fix a full-rank lattice $\Lambda\subset\R^2$ and a positive finite number of bounded
$C^\infty$ strictly convex obstacles $D_1,\ldots,D_s$ with everywhere
positive curvature. All distinct translates $D_a+\lambda$, $\lambda\in
\Lambda$, have disjoint closures. Let $A>0$ be the free area per lattice
cell. The unit-speed billiard is prepared with normalized phase volume
$\dd q\dd\vartheta/(2\pi A)$ on the torus complement. We observe the
collision count $N_t$ and, when specified, the first and last collision
positions. Finite horizon, congruence of the obstacles, and symmetry
are not assumed for the general forward and boundary-law results.

Consider a fixed such configuration and a smooth finite-dimensional
family $D_{a,\xi}$ through it. The parameter $\xi$ ranges over a compact
set $K_0$ in a sufficiently small neighborhood of the reference parameter.
The general geometric conclusions are local about an arbitrary reference configuration,
not only about circles. A compact family is covered by finitely many
such neighborhoods, with common constants after taking finite extrema.
Boundaries may be represented by smooth support functions. At every
fixed derivative order, constants can depend on the corresponding
higher norms of this family, on its positive separation and curvature
bounds, and on the lattice. They do not depend on the collision order.

Let $\mathcal E$ be the finite set of oriented closest pairs attaining
the minimum gap in the reference configuration, modulo common lattice
translation. An index $e=(a,b,\lambda)$ represents a segment from
$D_{a,\xi}$ to $D_{b,\xi}+\lambda$; reversing it gives
$\bar e=(b,a,-\lambda)$. Section~\ref{sec:g-localization} proves that
these pairs continue smoothly, remain unobstructed, and include every
channel contributing to a common ground-onset collar. Write
\begin{equation}\label{eq:g-general-data}
 \begin{split}
 g_*&=\min_{e\in\mathcal E}g_e,\qquad
 c_{e,b}=1+g_e\kappa_{e,b}\quad(b=0,1),\\
 c_e&=\sqrt{c_{e,0}c_{e,1}},\qquad \gamma_e=\arcosh c_e,
 \qquad d_{j,e}=t-jg_e.
 \end{split}
\end{equation}
Here $\kappa_{e,0}$ and $\kappa_{e,1}$ are the positive curvatures at
its two contact points. They need not be equal. A channel event
$E_{j,e}(d)$ consists of the $j+1$ impacts alternating between the
specified facing patches in the window $t=jg_e+d$. These patches are
fixed around the reference contacts and chosen before observation.
They identify a measurable part of a physical collision record, not an
extra latent observation. At a nonminimal channel's own onset this
selection is required; the unselected count law there is a different
observation.

Parameter derivatives are multi-index derivatives in a fixed Euclidean
parameter chart; offset derivatives hold $d$ fixed as an independent
coordinate. Matrix norms are Euclidean operator norms unless a trace
norm is explicitly indicated. The growing record arrays later use the
maximum norm. These conventions distinguish uniform normalized
amplitudes from physical-time derivatives, which can carry powers of $j$.


\subsection{The nonlinear relative law}
Fix an oriented channel $e$ and abbreviate $g=g_e$, $\gamma=\gamma_e$.
Use transverse endpoint coordinates $u,v$ at its two contacts. Let
$W_j(u,v)$ be the stationary length of its $j$ alternating flights,
$E_j=W_j-jg$, $d_j^0=-W_{j,uv}(0,0)$, and $b_j=-W_{j,uv}/d_j^0$.
For $p=j\bmod2$ put
\[
 a_b=\frac{\sqrt{c_0c_1}\sinh\gamma}{g c_{1-b}},\qquad
 q_p=\sqrt{a_0a_p},\qquad d_j^0=\frac{q_p}{\sinh(j\gamma)}.
\]
The stationary bridge and these formulas are proved in
Section~\ref{sec:g-action}. All norms in the next theorem are on fixed
common neighborhoods and include the indicated geometric derivatives.

\begin{theorem}[Relative boundary law]\label{thm:v8-main-relative}
There are positive smooth functions $B_b$ and strictly convex smooth
functions $S_b$, $b=0,1$, with
\[
 S_b(0)=S_b'(0)=0,\qquad S_b''(0)=a_b,\qquad B_b(0)=1,
\]
a common endpoint box, a number $d_*>0$, and $0<\tau<1$ such that,
for every fixed derivative order $k$ and every $j\ge1$,
\begin{equation}\label{eq:v8-relative}
 \|E_j-S_0\oplus S_p\|_{C^k}
 +\|b_j-B_0\otimes B_p\|_{C^k}\le C_k\tau^j,
                    \qquad p=j\bmod2.
\end{equation}
The action difference vanishes to second order at the origin. Define
\begin{equation}\label{eq:v8-boundary-integral}
 \mathcal F_p(d)=\frac{q_p}{\pi d^2}
  \int(d-S_0(u)-S_p(v))_+B_0(u)B_p(v)\,\dd u\dd v.
\end{equation}
Then $\mathcal F_p$ extends smoothly to $[0,d_*]$, is positive there,
and has value one at zero. The exact selected-channel probabilities obey
\begin{equation}\label{eq:v8-relative-integral}
 \left\|\frac{2A\sinh(j\gamma)}{d^2}
           \Pr(E_{j,e}(d))-\mathcal F_p(d)\right\|_{C^k([0,d_*])}
                                  \le C_k\tau^j.
\end{equation}
The norm includes every fixed mixed geometric and offset derivative.
The limiting conditional endpoint--residual law at $d>0$ has density
proportional to
\[
 B_0(u)B_p(v)\one_{\{r>0,\ S_0(u)+S_p(v)+r<d\}}.
\]
All assertions hold uniformly on the compact local families specified
above. For analytic families the endpoint constructions are analytic.
\end{theorem}

\begin{proof}[Proof and dependencies]
Lemma~\ref{lem:g-relative} constructs the finite bridge and its relative
determinant. Lemma~\ref{lem:v4-halfline} constructs $S_b,B_b$ with their
normalizations. Theorem~\ref{thm:v4-factorization} proves
\eqref{eq:v8-relative} by comparing two separated perturbation blocks in
trace norm, including their mixed derivatives. Theorem~\ref{thm:v4-law}
proves \eqref{eq:v8-relative-integral} using the common Morse domain of
Lemma~\ref{lem:g-radial}, and identifies the conditional density.
These proofs are supplied in full in the following sections.
\end{proof}

\subsection{Comparison with spectral rigidity}
\label{sec:v13-literature}
The closest comparisons differ both in their information sets and in
the geometry they recover. We record these differences at the level
of the stated theorems, rather than using the word ``spectral'' to
identify distinct observations.

De Simoi--Kaloshin--Leguil~\cite{DSKL}, Definition~1.3 and the Main
Theorem, consider analytic open dispersing billiards with at least
three obstacles and the non-eclipse condition. Their symmetric class
requires equal curvature jets at the two endpoints of a selected
period-two orbit, with both curvature jets even in arclength. On an
open dense nondegenerate subclass, the marked periodic-orbit lengths
determine the whole table up to isometry. Their reconstruction uses
hyperbolic Birkhoff and homoclinic data; the genericity condition includes
a nonzero first Birkhoff coefficient. Our even-contact inverse instead
supplies the gap and separate labelled curvatures and observes selected
near-onset probability germs. It recovers the participating contacts,
not an arbitrary whole table, while allowing the two contact jets to
vary independently. The common triangular and analytic-continuation
steps do not identify these information sets or their conclusions.

Finamore--Leguil~\cite{FinamoreLeguil}, Theorem~A in the cited version,
prove isometry for diffeomorphic finite-horizon Sinai billiards with
the same \emph{enriched} marked length spectrum. Their setting has
$C^r$ boundaries, $r\ge3$. The enrichment is defined through limiting
geodesic data on surfaces lifting the billiard and includes cycles
that are not ordinary billiard periodic orbits. Their mechanism uses
$\mathrm{CAT}(0)$ geometry and an adaptation of marked-length rigidity;
it is a global geometric theorem. Here the data are selected local
probabilities, the general target is a pair of smooth energy profiles,
and no finite horizon is assumed. The calibration theorem removes
$g,A,\gamma$, but retains channel labels and uniform class certificates.
It neither obtains the enriched spectrum nor derives that global
isometry conclusion from these local observations.

Zelditch~\cite{Zelditch}, Theorem~1.1, treats simply connected bounded
analytic plane domains in the class $\mathcal D_{1,L}$, with an
involution reversing a prescribed bouncing-ball orbit and the stated
nondegeneracy and length-multiplicity hypotheses. Either the Dirichlet
or Neumann Laplace spectrum determines the domain within that class;
Theorem~1.4 treats the specified dihedral classes. The proof extracts
highest boundary derivatives from wave-trace invariants by oscillatory
boundary integrals and stationary phase. This is an important comparison
for highest-jet extraction, but the Laplace spectrum is not an A2
observation. Our two-flight coefficient is an integrated physical flux,
including its varying mixed derivative, and uses supplied labelled
leading geometry. Our general smooth-profile theorem does not turn
symmetrized energy data into unrestricted analytic-domain rigidity.

The distinction also concerns what remains to be proved after elementary
identities are available. Schur complements and Jacobi Green functions
are classical, as is the relation between action Hessians and
monodromy~\cite{Hill}. Section~\ref{sec:v9-operators} isolates those
identities from the nonlinear, differentiated trace-norm comparison
and its transfer through physical integration. The full-profile inverse
is compared with nonlinear deautoconvolution in
Section~\ref{sec:v10-deautoconvolution}; the weighted Abel coordinate
is tied to the integrated physical flux and its observation topology.
Our principal assertion is the combination of uniform nonlinear relative
control on a nonshrinking collar, determination of the resulting
smooth invariant without symmetry, and a charged observation procedure
for that invariant. It is not a claim that the cited global inverse
problems have been subsumed, or an exhaustive priority assertion.

\subsection{Organization}
The first part proves geometric localization, the relative boundary law,
and nonlinear energy compatibility, followed by the two-contact limiting
inverse and its physical realization. The two-flight theorem then gives
the direct finite-record comparison and strengthens the observation
design to $j_0=2$. The second part proves Abel stability and the
regularized complete-profile observation, including self-calibration.
The appendix contains all auxiliary endpoint experiments,
finite-dimensional and smooth-remainder inverses, earlier acquisition
constructions, contact-jet proofs, and collision-record formulas.
These additional experiments are not used to replace the hypotheses
of the function-valued inverse.
